\section{Rational functions, multipliers, and block-separated certificates}
\label{sec:fields-certificates}

The obstructions of Section~\ref{sec:fields} concern polynomial squares.
Two changes of the certificate format are common in exact certification.
One allows denominators: a nonnegative rational polynomial is a sum of
squares of rational functions by Artin's theorem, and exact certificates with
rational-function or multiplier terms are an established strategy
\cite{KaltofenLiYangZhi2012}. The other imposes structure: sparse
sum-of-squares methods split a certificate into blocks of variables
\cite{KojimaKimWaki2005}. This section measures both changes on the
families of Section~\ref{sec:fields}.
\begin{enumerate}[label=(\alph*)]
\item A quadratic in the span of the square factors whose homogeneous
  quadratic part is positive definite,
  together with an affine ideal membership of the subtracted quadratics,
  gives an explicit rational certificate with the multiplier
  $\omega=1+\norm X^2$ (Theorem~\ref{thm:fields-multiplier}); higher-degree
  membership gives the multiplier $\omega^d$
  (Theorem~\ref{thm:fields-higher-membership}). After a polynomial-size
  increase of the scale, the tower quartics, whose polynomial certificates
  need a coefficient of degree $5^k$, have short rational-function
  certificates with an everywhere-positive quadratic common denominator, and
  two is the minimum denominator degree
  (Corollary~\ref{cor:fields-tower-denominator}).
\item A prescribed multiplier behaves differently. In a scaling family of
  ternary quartics with a fixed Hessian at the minimizer, no fixed power of
  $\omega$ suffices (Theorem~\ref{thm:fields-radial}); every power below a
  bound of order $\log L/\log\log L$ in the input length $L$ fails, although
  an adapted quadratic denominator gives a certificate of size $O(L)$
  (Theorem~\ref{thm:fields-radial-height}).
\item Block-separated rational certificates can fail to exist, can require
  the field $\Q(2^{1/5^k})$, or can exist but require $\Omega(k2^k)$ bits,
  while an unrestricted rational certificate is short; this persists when both
  blocks are strongly SOS-convex (Theorems~\ref{thm:fields-blocks}
  and~\ref{thm:fields-blocks-height}).
\end{enumerate}
These are statements about specified certificate formats. None is a lower
bound for unrestricted certificates or for the complexity of deciding
nonnegativity.

\subsection{Certificate formats}
\label{sec:fields-formats}

Let $f\in\Q[X]$. A \emph{rational-function sum of squares with common
denominator} $d\in\Q[X]\setminus\{0\}$ is an identity
$f=\sum_j(u_j/d)^2$ of rational functions with $u_j\in\Q[X]$; equivalently
$d^2f=\sum_ju_j^2$. It is \emph{regular} if $d$ has no real zero. A
\emph{multiplier certificate} is an identity $uf=\sum_jg_j^2$ with
$g_j\in\Q[X]$ and a fixed positive polynomial $u$. Put
\[
 \omega(X)=1+\norm X^2 ,
\]
and define the \emph{radial order} $\nu(f)\in\N\cup\{\infty\}$ as the least
$N\ge0$ with $\omega^Nf$ a sum of squares in $\Q[X]$. Since $\omega$ is a
rational sum of squares, $\omega^Nf$ is a rational sum of squares for every
$N\ge\nu(f)$. A multiplier is not a common denominator: from
$\omega f=\sum_jg_j^2$,
\begin{equation}
 f=\sum_j\Bigl(\frac{g_j}{\omega}\Bigr)^2
   +\sum_{i=1}^n\sum_j\Bigl(\frac{X_ig_j}{\omega}\Bigr)^2 ,
 \label{eq:fields-multiplier-to-denominator}
\end{equation}
a regular certificate with denominator $\omega$ of degree two and numerators
of degree $\max_j\deg g_j+1$; the cleared identity is
$\omega^2f=\omega\sum_jg_j^2$.

For disjoint blocks of variables $x$ and $y$ and a polynomial $f(x)+g(y)$, a
\emph{block-separated sum of squares over $E$} is a sum of squares over $E$
in which every square factor lies in $E[x]$ or in $E[y]$. A \emph{separated
Gram certificate over $E$} is a pair of positive semidefinite matrices
$Q_x,Q_y$ with entries in $E$, each on the monomial vector of its own block,
including its own constant monomial, with
$f(x)+g(y)=z(x)^{\mathsf T}Q_xz(x)+z(y)^{\mathsf T}Q_yz(y)$.

\subsection{Quadratic denominators from ideal membership}
\label{sec:fields-multipliers}

Denominators of degree one never help.

\begin{lemma}[Affine denominators cancel]
\label{lem:fields-affine-denominator}
If $f\in\Q[X]$ has a rational-function sum of squares with a common
denominator of degree at most one, then $f$ is a sum of squares in $\Q[X]$.
\end{lemma}

\begin{proof}
A constant denominator can be absorbed. Otherwise $d$ is a nonconstant
rational affine polynomial and $d^2f=\sum_ju_j^2$. On the real hyperplane
$d=0$ every $u_j$ vanishes. After a rational affine change of coordinates
making $d$ a multiple of $X_1$, the polynomial $u_j(0,X_2,\ldots,X_n)$
vanishes on $\R^{n-1}$, so it is zero and $X_1$ divides $u_j$ in $\Q[X]$.
Thus $d$ divides every $u_j$, and cancelling $d^2$ gives a polynomial sum of
squares.
\end{proof}

Consequently every rational-function certificate of a polynomial that is not a
rational sum of squares, such as $f_\ell$ or $f_{k,\lambda}$, has a common
denominator of degree at least two. The next theorem shows how degree two is
achieved, constructively and with a short certificate, whenever two pieces
of algebraic data are available.

\begin{theorem}[Quadratic multipliers from affine membership]
\label{thm:fields-multiplier}
Let $g_1,\ldots,g_m\in\Q[X_1,\ldots,X_n]$ have degree at most two, and put
$F_0=\sum_ig_i^2$. Suppose the input also supplies
\begin{enumerate}[label=(\roman*)]
\item rationals $\kappa_i$ such that
  $G=\sum_i\kappa_ig_i=X^{\mathsf T}HX+\ell_G^{\mathsf T}X+c_G$ has $H\succ0$;
\item rational quadratics $R_1,\ldots,R_s$ with affine rational multipliers,
  $R_a=\sum_i(\alpha_{ai}+\beta_{ai}^{\mathsf T}X)g_i$;
\item a rational symmetric $s\times s$ matrix $J$, of any signature.
\end{enumerate}
One can compute, in time polynomial in the input length, a positive integer
$\lambda_*$ with the following property. For every rational
$\lambda\ge\lambda_*$, the polynomial $\omega(\lambda F_0-R^{\mathsf T}JR)$
has a rational positive definite Gram matrix on an explicit vector of
polynomials of degree at most three, and hence is a sum of squares of
rational polynomials of degree at most three, with a certificate of size
polynomial in the input and the length of $\lambda$. If a rational positive
definite full Hessian Gram of $F_0$ is supplied as well, $\lambda_*$ can also
be chosen so that $\lambda F_0-R^{\mathsf T}JR$ has a rational full Hessian
Gram at least $I$ for all $\lambda\ge\lambda_*$. If the $g_i$ have a common
real zero $p$, then every $R_a$ vanishes at $p$, and so do the value and the
gradient of $\lambda F_0-R^{\mathsf T}JR$.
\end{theorem}

The proof in Appendix~\ref{app:fields-multiplier} rests on an exact identity
that represents zero by a matrix with a positive definite block. On the list
$w=(X_jg_i)_{0\le j\le n,\,1\le i\le m}$, with $X_0=1$, the identity matrix
represents $\omega F_0$. The new polynomials $X_jR_a$ carry the block
$I_s\otimes H$ in a Gram representation of $0$, whose correction terms are
fixed by the identity
$GR_a^2-GR_a(X^{\mathsf T}T_aX+r_a^{\mathsf T}X+\sigma_a)=0$, where
$R_a=X^{\mathsf T}T_aX+r_a^{\mathsf T}X+\sigma_a$. A small positive multiple of
this zero representation added to the identity matrix gives a positive
definite Gram matrix of $\omega F_0$ on the enlarged list, which then
dominates the Gram matrix of $\omega R^{\mathsf T}JR$ after scaling. No
convexity of $F_0$ and no common zero are needed.

Affine membership is a genuine restriction. In $\Q[x,y,z]$ the quadratics
$g_1=x^2+y^2+z^2-1$, $g_2=x^2$, $g_3=xy-z$, $g_4=yz-1$ have no common
complex zero, $g_1$ has positive definite quadratic part, and
\[
 1=xg_1+(-x+y^2-1)g_2+(-xy+xz-y)g_3+(-x^2-x-1)g_4 ,
\]
yet $1$ is not in the span of the $g_i$ and $X_jg_i$
(Appendix~\ref{app:fields-membership-example}). The following extension
handles membership of higher degree.

\begin{theorem}[Multipliers from higher-degree membership]
\label{thm:fields-higher-membership}
Let $g_i$, $F_0$, $G$, and $J$ be as in Theorem~\ref{thm:fields-multiplier},
let $d\ge1$, and let rational quadratics $R_a$ be supplied with
representations $R_a=\sum_iu_{ai}g_i$, $u_{ai}\in\Q[X]$, $\deg u_{ai}\le d$.
One can compute a positive integer $\lambda_*$ such that for every rational
$\lambda\ge\lambda_*$ the polynomial $\omega^d(\lambda F_0-R^{\mathsf T}JR)$
is a sum of squares of rational polynomials of degree at most $d+2$. Time and
output size are polynomial in the input length, $n$, $m$, $s$, $d$, and
$\binom{n+d}d$. The real zero set of $\lambda F_0-R^{\mathsf T}JR$ is the
common real zero set of the $g_i$. In particular this polynomial has a
regular rational-function sum of squares with common denominator
$\omega^{\lceil d/2\rceil}$.
\end{theorem}

The size bound counts the expanded basis of monomials of degree at most $d$;
it is not polynomial in a binary encoding of $d$. The targets must remain
quadratic: a negative square of higher degree could make
$\lambda F_0-R^{\mathsf T}JR$ negative at infinity for every $\lambda$. The
proof (Appendix~\ref{app:fields-higher}) attaches one zero identity to each
monomial multiple $X^\alpha R_a$ and weights the identities of level
$|\alpha|+1$ by $\rho^{2(|\alpha|+1)}$ for a small rational $\rho$; a single
diagonal congruence then shows that every correction term carries at least
one factor $\rho$. The example above shows that the class with $d=2$ is
strictly larger than the affine class; it does not show that the multiplier
$\omega^2$ is necessary there.

\begin{corollary}[Quadratic denominators for large least fields]
\label{cor:fields-tower-denominator}
\leavevmode
\begin{enumerate}[label=(\alph*)]
\item For every $k\ge1$ one can compute, in time polynomial in $k$, an
  integer $\lambda'_k\ge\lambda_k$ such that for every rational
  $\lambda\ge\lambda'_k$ the tower quartic $f_{k,\lambda}$ of
  Theorem~\ref{thm:fields-tower} satisfies
  $\omega f_{k,\lambda}\in\Sigma\Q[X]^2$, with square factors of degree at
  most three and a certificate of size polynomial in $k$ and the length of
  $\lambda$. Hence $f_{k,\lambda}$ has a regular rational-function sum of
  squares with common denominator $\omega$, and the minimum degree of a common
  denominator of a rational-function sum of squares of $f_{k,\lambda}$ is two.
  All conclusions of Theorems~\ref{thm:fields-tower}
  and~\ref{thm:fields-individual} hold for $f_{k,\lambda}$.
\item For every prime $\ell\ge5$ one can compute, in time polynomial in
  $\ell$, an integer $c_\ell\ge1$ such that for every integer $c\ge c_\ell$
  the integer quartic
  $f_\ell^{(c)}=c(f_\ell+r^2)-r^2$ of Remark~\ref{rem:fields-prime-scales}
  satisfies $\omega f_\ell^{(c)}\in\Sigma\Q[X]^2$. The monomial count in
  Theorem~\ref{thm:fields-prime}(a) and its conclusions (b)--(d) hold for
  $f_\ell^{(c)}$. Given $c$, construction takes time polynomial in $\ell$
  and the binary length of $c$; the coefficient and Hessian-Gram entry
  lengths are $O(\log\ell+\log c)$.
\end{enumerate}
\end{corollary}

\begin{proof}
(a) By construction (Appendix~\ref{app:fields-tower}),
$F_{k,0}=(G_k/(t\nu))^2+\sum_{i,j}(r_{i,j}/\nu)^2$, $j\in\{1,2,4\}$, for
positive rationals $t,\nu$. Apply Theorem~\ref{thm:fields-multiplier} to
these factors, so that its $F_0$ is exactly $F_{k,0}$, with $\kappa$
selecting the factor $G_k/(t\nu)$, whose quadratic part is positive
definite, and the single target $R_1=r_{k,3}$ with $J=(1)$. The exact
identity
\[
 r_{k,3}=x_kr_{k,2}-y_kr_{k,1}
 =\nu x_k\,(r_{k,2}/\nu)-\nu y_k\,(r_{k,1}/\nu)
\]
is an affine membership in these factors. The theorem gives $\lambda_*$ with
$\omega(\lambda F_{k,0}-r_{k,3}^2)\in\Sigma\Q[X]^2$ for every rational
$\lambda\ge\lambda_*$; take $\lambda'_k=\max\{\lambda_k,\lambda_*\}$. The
field proof of Theorem~\ref{thm:fields-tower} holds for every
$\lambda\ge\lambda_k$, and $f_{k,\lambda}$ is not a rational sum of squares,
so Lemma~\ref{lem:fields-affine-denominator} excludes denominators of degree
at most one; \eqref{eq:fields-multiplier-to-denominator} gives degree two.
(b) Apply the theorem to the factors $A$ and $(N_1/N_0)C_i$ of
$f_\ell+r^2=A^2+(N_1/N_0)^2\sum_iC_i^2$ (Appendix~\ref{app:fields-prime}),
with $\kappa$ selecting $A$, whose quadratic part is positive definite, and
$J=(1)$. The identity
\[
 r=x_0C_1-x_1C_0=\frac{N_0}{N_1}\Bigl(x_0\,\frac{N_1}{N_0}C_1
   -x_1\,\frac{N_1}{N_0}C_0\Bigr)
\]
is an affine membership with rational multipliers, and
$c(f_\ell+r^2)-r^2=f_\ell^{(c)}$; take $c_\ell=\lambda_*$.
\end{proof}

For $\lambda\ge\lambda'_k$, every polynomial sum of squares of
$f_{k,\lambda}$ with algebraic coefficients needs one coefficient of degree at least
$5^k$, while a rational-function sum of squares with the fixed denominator
$\omega$ has polynomial size. Without any size bound, some power of $\omega$
works for each such quartic, with an exponent that depends on the quartic.

\begin{proposition}[Finite radial order]
\label{prop:fields-radial-finite}
Let $f\in\Q[X_1,\ldots,X_n]$ be nonnegative on $\R^n$, of even degree
$2e\ge2$, with positive definite homogeneous part of degree $2e$, finitely
many real zeros, and positive definite Hessian at each real zero. Then
$\omega^Nf$ is a sum of squares in $\Q[X]$ for some integer $N\ge0$, that
is, $\nu(f)<\infty$. In particular every certified quartic with minimum zero
has finite radial order.
\end{proposition}

Nonnegativity cannot be omitted: $x^2-1$ satisfies the other conditions, but
$\omega^N(x^2-1)$ is negative at the origin. A certified quartic with minimum
zero satisfies all the hypotheses: its quartic form is positive definite
(proof of Corollary~\ref{cor:fields-dimension}), its only real zero is the
minimizer, and its Hessian is positive definite everywhere. The proof
(Appendix~\ref{app:fields-radial-finite}) works in the coordinate ring of the
unit sphere, where sums of squares form an archimedean preordering, and
applies the pure-state and order-unit theory of Burgdorf, Scheiderer, and
Schweighofer \cite{BurgdorfScheidererSchweighofer2012} to the square of the
rational vanishing ideal of the zeros of the homogenized polynomial. Rational
coefficients come from working in a ring over $\Q$ throughout; a statement
about real coefficients alone would not give them. The proposition gives no
bound on $N$ and no certificate size, and the lower bounds of
Section~\ref{sec:fields-radial} do not use it.

\begin{example}[The explicit ternary quartic]
\label{ex:fields-ternary-multiplier}
For the quartic $F_\star$ of Example~\ref{ex:fields-ternary}, which is not a
rational sum of squares, Appendix~\ref{app:fields-ternary-multiplier} prints
a vector $b$ of fifteen rational polynomials of degree at most three, which is
a basis of the rational cubics vanishing at $p_\star$, and an integer
symmetric matrix $N$ with entries of absolute value at most $4240$, such that
\[
 b^{\mathsf T}Nb=8\,\omega F_\star,\qquad N-8I\succ0,
\]
the latter certified by fifteen positive leading principal minors. Thus
$\omega F_\star$ is a rational sum of squares of cubics, $F_\star$ has a
regular rational-function sum of squares with denominator $\omega$, and two is
the minimum common-denominator degree. Here the scale of $F_\star$ is not
increased; the certificate is a finite exact identity.
\end{example}

\subsection{Prescribed radial multipliers}
\label{sec:fields-radial}

The denominator in Corollary~\ref{cor:fields-tower-denominator} is the fixed
polynomial $\omega$, but the scale of the quartic was chosen to fit it. If
instead we prescribe the radial hierarchy $\omega^N$ for every quartic in a
scaling family, no fixed order $N$ suffices. Let $f^\circ=f_5^{(c_5)}$ be
the ternary integer quartic of Corollary~\ref{cor:fields-tower-denominator}(b)
in the variables $(x_0,x_1,x_2)$. It has a rational positive definite full
Hessian Gram, $\nabla^2f^\circ\succeq I$, minimum zero at
$p=(a^2,a^3,a^4)$ with $a=2^{1/5}$, $\varphi_5(f^\circ)=-4$ for the
functional $\varphi_5=[x_2]+[x_0^2]$ of
\eqref{eq:fields-prime-functional}, and $\omega f^\circ\in\Sigma\Q[X]^2$. For
integers $t\ge1$ put
\[
 f_t(X)=t^{-2}f^\circ(tX).
\]

\begin{theorem}[Unbounded radial order]
\label{thm:fields-radial}
\leavevmode
\begin{enumerate}[label=(\alph*)]
\item Each $f_t$ has a rational positive definite full Hessian Gram,
  $\nabla^2f_t\succeq I$, unique minimizer $p/t$ with value zero, and
  $\nabla^2f_t(p/t)=\nabla^2f^\circ(p)$. Its binary input length satisfies
  $L(t)=\Theta(\log t)$.
\item $(1+t^2\norm X^2)f_t$ is a sum of squares of rational cubics with a
  certificate of size $O(L(t))$. Hence $f_t$ has a regular rational-function
  sum of squares with a quadratic common denominator, and two is the minimum
  degree of a common denominator.
\item For every $N\ge0$ one can compute an integer $\tau_N$ such that
  $\omega^Nf_t$ is not a sum of squares in $\Q[X]$ for every integer
  $t\ge\tau_N$. More generally, for each fixed $u\in\Q[X]$ with $u(0)>0$, the
  polynomial $uf_t$ is not a sum of squares in $\Q[X]$ for all sufficiently
  large integers $t$.
\end{enumerate}
\end{theorem}

\begin{theorem}[Quantitative radial order]
\label{thm:fields-radial-height}
There are constants $c>0$ and $t_0$, depending only on $f^\circ$, such that
\[
 \nu(f_t)\ge c\,\frac{\log\log t}{\log\log\log t}
 \qquad\text{for every integer }t\ge t_0 .
\]
Equivalently $\nu(f_t)=\Omega(\log L/\log\log L)$ with $L=L(t)$, while
$f_t$ has an adapted rational certificate of size $O(L)$. Each $\nu(f_t)$ is
finite by Proposition~\ref{prop:fields-radial-finite}.
\end{theorem}

The proofs are in Appendix~\ref{app:fields-radial}. The mechanism is a
sequence of rational linear functionals $\psi_d$ on polynomials of degree at
most $2d$ that are positive on squares of real combinations of rational
polynomials of degree at most $d$ vanishing at $p$, and negative on
$f^\circ$. The first is $\varphi_5$ plus a small multiple of the evaluation
sum over the grid $\{0,1,2\}^3$; each further one adds a large integer
multiple of a grid functional of top degree, chosen through an exact Schur
complement. A rational sum of squares of $\omega^Nf_t$ would, after the
substitution $X=x/t$, have factors in the domain of $\psi_{N+2}$, and
$\psi_{N+2}$ is negative on it once $t$ is large. Tracking the heights of the
functionals gives $\log\log\tau_N=O(N\log N)$, which inverts to
Theorem~\ref{thm:fields-radial-height}.

The scaling-and-closedness method is Reznick's
\cite[Theorem~1 and Corollary~2]{Reznick2005}, who used it to show that no
single multiplier, and no finite menu of multipliers, works for all
nonnegative forms of a given degree and number of variables outside the
cases of Hilbert's theorem. Here every $f_t$ is a real sum of squares
already at order zero and is strongly SOS-convex with a strict rational
certificate; the obstruction
is purely arithmetic and comes from the rational vanishing space of the
irrational minimizer. The Hessian at the minimizer is the same for all $t$,
so the growth of the order is not explained by local ill-conditioning. The
theorem concerns one prescribed hierarchy. It does not say that real
sum-of-squares optimization needs high order, and it does not say that finding
an adapted denominator is hard; here the scale makes it explicit.

\subsection{Block-separated certificates}
\label{sec:fields-blocks}

Sparse certificates restrict the variables that may appear together in a
square. For a separable polynomial $f(x)+g(y)$, the strongest such
restriction asks for squares in $x$ alone or $y$ alone. Real sparse
decompositions of this kind are studied in \cite{KojimaKimWaki2005}, and for
the separable quartics below real separated certificates exist
(Theorem~\ref{thm:fields-blocks}(b) with $E=\R$). Over $\Q$ the answer
differs. The constructions use a lift of a certified quartic to a quadratic
graph.

\begin{lemma}[A rational quadratic lift]
\label{lem:fields-lift}
Let $(f,A)$ be a certified quartic in $n$ variables with minimum zero at $p$;
the point $p$ is not part of the input. Introduce variables $y\in\R^n$ and
$z=(z_{ij})_{1\le i\le j\le n}$. In time polynomial in the size of $(f,A)$ one
can construct a rational quartic $g(y,z)$ and a rational sum of squares of
$f(x)+g(y,z)$ with quadratic factors, of polynomial size, such that $g\ge0$
and $g(p,(p_ip_j)_{i\le j})=0$.
\end{lemma}

The construction (Appendix~\ref{app:fields-lift}) replaces a quadratic
pair in each cubic monomial of $\nabla f(y)$ by a variable $z_{ij}$ to obtain
a quadratic map $a(y,z)$, and sets
$g=\norm{a}^2/(2\mu)-f(y)+\lambda\sum_{i\le j}(y_iy_j-z_{ij})^2$ with
$\mu=\mu_A/2$ and an explicit integer $\lambda$. The joint certificate
combines the Taylor Gram of Lemma~\ref{lem:taylor-sos} in the difference
$x-y$ with a completed square. Substituting $x=p$ shows $g\ge0$.

\begin{theorem}[Rational block separation fails]
\label{thm:fields-blocks}
Let $k\ge1$, $f=f_{k,\lambda_k}$, and let $g$ be its lift from
Lemma~\ref{lem:fields-lift}, in the variables $w=(y,z)$. Let $B_k(w)$ be the
rational quartic obtained by applying Theorem~\ref{thm:quadratic-graph} to
$f$ and its Hessian Gram, with full Hessian Gram $A_{B_k}\succ0$, and let
$C$ be a rational symmetric full Hessian Gram of $g$. Put
\[
 T_k=\Bigl\lceil\frac{1+\sum_{i,j}|C_{ij}|}{\mu_{A_{B_k}}}\Bigr\rceil,
 \qquad \tilde g=g+T_kB_k .
\]
\begin{enumerate}[label=(\alph*)]
\item The quartic $f(x)+\tilde g(w)$ has a rational sum of squares with
  quadratic factors, of size polynomial in $k$. Both $f$ and $\tilde g$ have
  rational full Hessian Grams at least $I$, $\tilde g\ge0$, and the unique
  zero of $\tilde g$ is $w_*=(p_k,(p_ip_j)_{i\le j})$.
\item For every subfield $E\subseteq\R$ the following are equivalent:
  $2^{1/5^k}\in E$; $f(x)+\tilde g(w)$ has a block-separated sum of squares
  over $E$; it has a separated Gram certificate over $E$. In particular, for
  $k=1$ no rational block-separated certificate exists, although the joint
  rational sum of squares has fixed size and real separated certificates
  exist.
\item Every separated certificate with algebraic coefficients has one
  coefficient of degree at least $5^k$. The total number of variables is
  $3k+3k+\frac{3k(3k+1)}2=\Theta(k^2)$, so this degree is
  $5^{\Theta(\sqrt{n})}$ in the total dimension $n$.
\item Every positive semidefinite full Hessian Gram matrix of
  $f(x)+\tilde g(w)$ is singular.
\end{enumerate}
\end{theorem}

\begin{proof}
(a) The rational full Hessian Gram $C+T_kA_{B_k}$ of $\tilde g$ is at least
$(T_k\mu_{A_{B_k}}-\sum_{i,j}|C_{ij}|)I\succeq I$. Since $g\ge0$, $B_k\ge0$,
$g(w_*)=0$, and $w_*$ is the only zero of $B_k$, the only zero of
$\tilde g$ is $w_*$. Adding $T_k$ times the rational sum of squares of $B_k$
to the certificate of Lemma~\ref{lem:fields-lift} gives the joint
certificate. (b) Let a block-separated certificate over $E$ be given, and
group its terms: $f(x)+\tilde g(w)=S_x(x)+S_w(w)$, where $S_x$ and $S_w$ are
sums of squares, or Gram forms with positive semidefinite matrices, over $E$.
Then $S_x-f=\tilde g-S_w$ depends on neither block, so it is a constant
$c\in E$. Evaluating $S_x=f+c$ at $p_k$ gives $c\ge0$, and evaluating
$S_w=\tilde g-c$ at $w_*$ gives $c\le0$. Thus $f$ itself has a certificate of
the same type over $E$, and Theorem~\ref{thm:fields-tower}(b) gives
$2^{1/5^k}\in E$. Conversely, if $2^{1/5^k}\in E$, then $f$ is a sum of squares
over $E$, and specializing the joint rational certificate at
$x=p_k\in E^{3k}$ shows that $\tilde g$ is too. For $k=1$ and $E=\R$ this gives
real separated certificates. (c) follows from (b) and
Theorem~\ref{thm:fields-individual}. (d) Let $x_j$ be a variable of the first
block and $v_i$ a direction coordinate of the second. The Hessian of
$f(x)+\tilde g(w)$ is block diagonal, and the second block depends only on
$w$, so the Hessian biform has no monomial $x_j^2v_i^2$. In the full joint
Hessian basis this monomial arises only as the square of the coordinate
$x_jv_i$. Hence every Gram matrix has a zero diagonal entry there, and a
positive semidefinite one has a zero row.
\end{proof}

The joint polynomial is strongly convex with a short rational Hessian sum of
squares, obtained by adding the two block certificates, but by (d) it has no
positive definite full Hessian Gram; strong convexity on each block does not
couple the blocks. The obstruction is about rational coefficients: each
block is a certified quartic with minimum zero, hence a real sum of squares
by Lemma~\ref{lem:taylor-sos}(e), so real block-separated certificates exist.
When the minimum is positive, rational separated certificates exist, but
they can be long.

\begin{theorem}[Exponentially long rational block certificates]
\label{thm:fields-blocks-height}
Let $f=f_{1,\lambda_1}$ and $\tilde g$ be the quartics of
Theorem~\ref{thm:fields-blocks} with $k=1$. For $k\ge1$ let $F_k^\circ$ and
$A_k^\circ$ denote the quartic $F_k$ and its rational positive definite full
Hessian Gram $A_k$ from the construction preceding
Theorem~\ref{thm:heights-witness}, with $M_k=1000^{k+3}$, and let $u_k$ be
the rational affine function defined there. Put
$h_k=\lceil3/\mu_{A_k^\circ}\rceil F_k^\circ+u_k^2$, a quartic whose $2k$
variables we rename $\xi$. Consider
\[
 P_k(x,w,\xi)=f(x)+\tilde g(w)+h_k(\xi)
\]
with the two blocks $x$ and $(w,\xi)$.
\begin{enumerate}[label=(\alph*)]
\item $P_k$ is a rational quartic in $12+2k$ variables with a rational sum of
  squares and a rational sum-of-squares Hessian certificate, both of size
  polynomial in $k$, and $\nabla^2P_k\succeq I$. Its minimum
  $m_k=\min h_k$ satisfies $0<m_k\le M_k^{-2^{k+1}}$.
\item Rational block-separated sums of squares of $P_k$ exist.
\item In every separated Gram certificate $(Q_x,Q_{w,\xi})$ over $\Q$, the
  constant entry of $Q_x$ has reduced denominator at least
  $M_k^{2^{k+1}}/b_0$, where $b_0$ is the denominator of $f(0)$. In every
  rational block-separated sum of squares, the constant terms of the square
  factors in $x$ have reduced denominators whose binary lengths sum to at
  least $2^k\log_2M_k-\frac12\log_2b_0$.
\end{enumerate}
In particular every rational separated certificate has $\Omega(k2^k)$ bits,
while the unrestricted rational certificate has size polynomial in $k$.
\end{theorem}

The proof is in Appendix~\ref{app:fields-blocks}. Its core is the same
grouping argument: a separated certificate determines a constant $c$ with
$S_x=f+c$, and now $0<c\le m_k$, because $c=0$ would certify $f$ over $\Q$ and
$c$ cannot exceed the minimum of the other block. A positive rational below
$M_k^{-2^{k+1}}$ has a long denominator, and the constant entry of the
$x$-block Gram is $f(0)+c$. The bound concerns two separately stored local
Gram matrices, each with its own constant monomial, or an unweighted
separated sum of squares; a single joint matrix with one shared constant
coordinate can hide the split. The lower bound is exponential in the
dimension and superpolynomial in the polynomially bounded input length; it is
not claimed to be $2^{\Omega(L)}$. Real separation questions for sparse
certificates are classical \cite{KojimaKimWaki2005}; the point here is that
a restriction harmless for real feasibility can remove all rational
certificates, or make them exponentially long, even under certified strong
convexity.
